\section{Bounded Support Stopping}

We keep the first gain edits and the second-slot eligibility
fixed. A credit becomes available at the first qualifying easy
query, expires after the next two queries, and is never forced
to release at the end of a recording. The continuation head uses
only current features, query ordinal and credit age. Its
training target compares immediate risk-adjusted utility
against the best remaining training opportunity or abstention.
Teachers exclude both donor and evaluation sites.
The learned continuation policy delays 37 of 293 support
edits. Against immediate release in the same recordings at
matched usage and amplitude it gives \StopSupportMatched{} in
conservative error reduction, and \StopSupportCV{} versus CV.
Neither contrast establishes an improvement, and the full risk
gate is unmet. The immediate policy without usage matching
passes \StopSupportImmediateJoint{} views; continuation and its
matched replay pass \StopSupportContinueJoint{} and
\StopSupportReplayJoint{}, respectively.

\begin{table}[h]
\centering
\caption{Outcome-aware capacity out of 606 views. Every column
keeps the first edits, second-slot recordings, second edit count
and amplitude fixed. Counts are not causal policy performance.}
\small
\begin{tabular}{lrrr}
\toprule
Queries & Top easy & All pred. easy & All legal \\
\midrule
Fixed release & \StopSupportFixedTop{} & \StopSupportFixedEasy{} & \StopSupportFixedLegal{} \\
Release through +2 & \StopSupportTwoTop{} & \StopSupportTwoEasy{} & \StopSupportTwoLegal{} \\
Release or later & \StopSupportLaterTop{} & \StopSupportLaterEasy{} & \StopSupportLaterLegal{} \\
Any query & \StopSupportAnyTop{} & \StopSupportAnyEasy{} & \StopSupportAnyLegal{} \\
\bottomrule
\end{tabular}
\end{table}

Expanding the recipient menu within the frozen easy threshold
changes little. Allowing all legal recipients raises capacity
even without delaying a query. This suggests testing admission
scores separately from waiting on these exposed development sites.
